\documentclass[conference]{IEEEtran}

\usepackage{booktabs}
\usepackage{amsmath}
\usepackage{graphicx}
\usepackage{url}

\title{ECE-9413 Final Report: Negacyclic NTT and SumCheck Prover}

\author{\IEEEauthorblockN{Jack}
\IEEEauthorblockA{New York University\\
ECE-9413}}

\begin{document}
\maketitle

\begin{abstract}
This report describes two exact finite-field kernels implemented in JAX: a
negacyclic number theoretic transform (NTT) and a SumCheck prover. Both
submitted implementations preserve the provided public APIs and pass the
required correctness tests. The NTT implementation uses a twisted
Cooley--Tukey schedule with Montgomery arithmetic and a GPU-oriented Pallas
path, while the SumCheck implementation uses a general expression evaluator,
serialized per-round folding, overflow-safe 32-bit arithmetic, and an optional
64-bit core track. On the local CPU backend, the NTT reached 0.633 ms median
latency for batch 4 and $N=16384$, or 103.59 Mcoeff/s. The required 32-bit
SumCheck prover reached about 1.05 ms for the linear expression at vars20 and
about 8.0 ms for the cubic expression. The main remaining limitation is that
these measurements are local CPU measurements; GPU speedups should be added
after running the same commands on the course GPU.
\end{abstract}

\section{Introduction}
NTT and SumCheck are both central kernels in modern cryptographic proof and
polynomial computation systems. They differ from ordinary dense floating-point
GPU workloads because all arithmetic is exact modular integer arithmetic, where
overflow behavior and reduction cost are part of the algorithm. The NTT is used
to accelerate polynomial multiplication over finite fields. SumCheck reduces a
large Boolean-hypercube sum to a sequence of low-degree univariate polynomial
claims, making it a key primitive in proof systems.

The project contributions are: (1) a batched negacyclic NTT using a
twist-plus-cyclic NTT schedule and Montgomery reduction, (2) a required 32-bit
SumCheck prover for the base expressions, (3) overflow-safe modular arithmetic
for both required 32-bit and optional 64-bit SumCheck tracks, and (4) benchmark
evidence tying correctness and performance claims to the submitted code. All
reported commands were run on the local JAX CPU backend on macOS. No GPU result
is claimed in this draft.

\section{Assignment 1: Negacyclic NTT}
\subsection{Design and Optimization}
The implemented NTT evaluates
\begin{equation}
y[k] = \sum_{n=0}^{N-1} x[n]\psi^{(2k+1)n} \pmod q.
\end{equation}
The code uses the standard negacyclic decomposition: first multiply each input
coefficient by the appropriate power of $\psi$, then run a cyclic
Cooley--Tukey NTT. This structure fits the provided tables and keeps the hot
path regular across the batch dimension.

The modular arithmetic hot path uses unsigned 32-bit values, with widening only
where multiplication requires it. The optimized path converts values into
Montgomery form in \texttt{prepare\_tables}. The precomputed tables include a
pretwisted $\psi$ table and Montgomery-form stage twiddles. The CPU path uses
JAX reshapes and vectorized butterfly stages. The GPU-oriented path builds a
Pallas kernel intended to fuse the stages for one row and then vmaps across the
batch. This keeps the public API unchanged while moving table preparation out
of the timed benchmark region.

The most useful optimization was moving expensive table conversion and
Montgomery constants into \texttt{prepare\_tables}. A less attractive approach
would have been direct evaluation, which is simpler but has $O(N^2)$ work and
does not scale. The final implementation stays with $O(N\log N)$ butterflies.

\subsection{Correctness}
Table~\ref{tab:a1-correctness} lists the correctness commands required by the
guidelines. The benchmark test wrapper must be run with \texttt{uv --directory
assignment1} from the repository root so its internal subprocess starts in the
assignment directory.

\begin{table}[h]
\caption{Assignment 1 correctness checks}
\label{tab:a1-correctness}
\centering
\begin{tabular}{ll}
\toprule
Command & Result \\
\midrule
\texttt{uv run pytest} via submission script & 5 passed \\
\texttt{pytest --logn 10 --batch 4} & 5 passed \\
\texttt{python -m tests.benchmark --tests --logn 10 --batch 4} & 5 passed \\
\bottomrule
\end{tabular}
\end{table}

\subsection{Performance}
Table~\ref{tab:a1-perf} reports a CPU sweep at batch 4. The best throughput in
this sweep occurs at $N=16384$, where the larger transform has enough work to
amortize fixed overhead.

\begin{table}[h]
\caption{Assignment 1 CPU benchmark sweep, batch 4}
\label{tab:a1-perf}
\centering
\begin{tabular}{rrrrr}
\toprule
$\log_2(N)$ & $N$ & Compile ms & Median ms & Mcoeff/s \\
\midrule
10 & 1024  & 462.84 & 0.057 & 71.29 \\
12 & 4096  & 144.06 & 0.293 & 55.88 \\
14 & 16384 & 161.24 & 0.633 & 103.59 \\
\bottomrule
\end{tabular}
\end{table}

\section{Assignment 2: SumCheck}
\subsection{Design and Optimization}
The SumCheck prover supports expressions represented as lists of additive
terms, where each term is a list of multiplicative factors. This general
evaluator supports the four required base expressions without hard-coding a
separate kernel per expression. For an expression of degree $d$, each round
emits $d+1$ evaluations of the round polynomial. Thus the linear expression
\texttt{a} emits two values per round, \texttt{a*b} and \texttt{a*b + c} emit
three, and \texttt{a*b*c} emits four.

The implementation strictly follows the required round order. For each round,
it first computes the round polynomial evaluations using the current tables.
Only after that round does it fold the tables using the provided challenge.
The code does not precompute challenge-folded tables ahead of time.

For the required 32-bit track, multiplication still widens to 64 bits because a
32-bit product can exceed 32 bits. Addition and subtraction, however, now avoid
unnecessary 64-bit casts. \texttt{mod\_add\_32} detects uint32 wraparound and
corrects by adding $2^{32}\bmod q$ before the final conditional subtraction.
\texttt{mod\_sub\_32} computes the wrapped case as $q-(b-a)$, avoiding
\texttt{a+q} overflow. This optimization keeps correctness on the edge cases
while reducing the integer width used by common MLE and expression operations.

The optional 64-bit core required a different approach. Since JAX does not
provide a native uint128 type, direct 64-bit multiplication can lose high
product bits. The submitted optional path uses two-limb Montgomery
multiplication with 32-bit limbs for \texttt{mod\_mul\_64}, and pairwise
overflow-safe modular addition for 64-bit reductions. This made the optional
core64 correctness suite pass, although the required report baseline remains
the 32-bit track.

\subsection{READMEs}
The Assignment 2 README was helpful for understanding the expected API, the
required expressions, and the exact correctness and benchmark commands. The
\texttt{sumcheck\_intro.md} walkthrough was especially useful for checking the
round order and confirming that challenge-based MLE updates must occur
between rounds rather than before the protocol begins. The custom-cases
documentation was useful for understanding possible extra experiments, but the
final required results here use the provided public cases.

\subsection{Correctness}
Table~\ref{tab:a2-correctness} reports the required 32-bit correctness
commands. Each command passed all five public cases and all four base
expressions for the selected variable count.

\begin{table}[h]
\caption{Assignment 2 required 32-bit correctness checks}
\label{tab:a2-correctness}
\centering
\begin{tabular}{ll}
\toprule
Command & Result \\
\midrule
\texttt{pytest --bits 32 --num-vars 4} & 20/20 cases passed \\
\texttt{pytest --bits 32 --num-vars 16} & 20/20 cases passed \\
\texttt{pytest --bits 32 --num-vars 20} & 20/20 cases passed \\
\texttt{bash scripts/make\_submission.sh} & 91 passed \\
\bottomrule
\end{tabular}
\end{table}

\subsection{Performance}
The required benchmark commands were run with \texttt{--runs 8 --warmup 3} on
the local CPU backend. Table~\ref{tab:a2-vars20} summarizes the final vars20
run. The latency ordering follows expression degree and operation count:
\texttt{a} is fastest, then \texttt{a*b}, then the two-term quadratic
\texttt{a*b + c}, then the cubic expression \texttt{a*b*c}. As degree grows,
the prover must emit more evaluation points per round and perform more modular
multiplications per table element.

\begin{table}[h]
\caption{Assignment 2 vars20 CPU benchmark, 32-bit}
\label{tab:a2-vars20}
\centering
\begin{tabular}{lrrr}
\toprule
Expression & Compile ms & Median ms & Mpts/s \\
\midrule
\texttt{a}       & 248--307 & 1.05--1.09 & 962--1002 \\
\texttt{a*b}     & 524--542 & 3.47--3.70 & 284--302 \\
\texttt{a*b + c} & 740--766 & 6.01--6.41 & 164--174 \\
\texttt{a*b*c}   & 943--979 & 7.98--8.45 & 124--131 \\
\bottomrule
\end{tabular}
\end{table}

Table~\ref{tab:a2-scaling} shows representative median latency as the table
size increases. The vars4 results are dominated by overhead, while vars16 and
vars20 expose the cost of scanning and folding much larger arrays.

\begin{table}[h]
\caption{Assignment 2 scaling by variable count, representative CPU medians}
\label{tab:a2-scaling}
\centering
\begin{tabular}{lrrr}
\toprule
Expression & vars4 ms & vars16 ms & vars20 ms \\
\midrule
\texttt{a}       & 0.005--0.006 & 0.13--0.16 & 1.05--1.09 \\
\texttt{a*b}     & 0.007--0.011 & 0.36--0.39 & 3.47--3.70 \\
\texttt{a*b + c} & 0.007--0.008 & 0.45--0.54 & 6.01--6.41 \\
\texttt{a*b*c}   & 0.007--0.008 & 0.54--0.59 & 7.98--8.45 \\
\bottomrule
\end{tabular}
\end{table}

\section{Cross-Assignment Discussion}
The main shared lesson is that modular arithmetic choices dominate both
assignments. In the NTT, Montgomery form is effective because each butterfly
does many modular multiplications with precomputed twiddles. In SumCheck, the
required 32-bit path benefits more from carefully avoiding unnecessary widened
addition and subtraction while still widening multiplication. The NTT is more
structured: every stage has a regular butterfly layout and a predictable
twiddle access pattern. SumCheck is more expression dependent: degree and term
count directly change the amount of work per round.

Memory movement also differs. NTT repeatedly permutes and combines arrays, so
the schedule and table layout matter. SumCheck repeatedly halves the tables by
MLE update; the largest rounds dominate runtime because later rounds quickly
become small. This also explains why vars4 benchmark numbers are mostly
overhead, while vars20 makes the differences between expressions visible.

\section{AI Use Reflection}
AI assistance was useful for quickly producing first-pass implementations,
debugging edge cases, and designing benchmark commands. The most useful
suggestion was to verify the optional 64-bit SumCheck path separately from the
required 32-bit path; this exposed a broken 64-bit multiplication routine even
though the required submission already passed. A plausible but poor suggestion
was to implement 64-bit multiplication with a simple shift-and-add loop. That
was correct but too slow for full vars16 and vars20 tests, so it was replaced
with a two-limb Montgomery method. The best workflow was to keep AI suggestions
tied to public tests and benchmark output rather than trusting code structure
alone.

\section{Conclusion}
Both assignments pass their required public tests with the submitted code.
Assignment 1 uses a Montgomery Cooley--Tukey negacyclic NTT and reached 0.633
ms median latency for batch 4 and $N=16384$ on the local CPU backend.
Assignment 2 uses a general SumCheck prover with serialized round folding and
overflow-safe modular arithmetic; its vars20 CPU medians range from about 1.05
ms for \texttt{a} to about 8.0 ms for \texttt{a*b*c}. The most important NTT
optimization was moving Montgomery conversion and table preparation outside the
hot path. The most important SumCheck optimization was keeping 32-bit add/sub
operations in uint32 while still handling overflow correctly. The main
remaining limitation is that this draft contains local CPU measurements only;
course GPU measurements should be added before final submission if GPU results
are expected.

\end{document}
